% Options for packages loaded elsewhere
\PassOptionsToPackage{unicode}{hyperref}
\PassOptionsToPackage{hyphens}{url}
%
\documentclass[
  11pt,
]{article}
\usepackage{amsmath,amssymb}
\usepackage{iftex}
\ifPDFTeX
  \usepackage[T1]{fontenc}
  \usepackage[utf8]{inputenc}
  \usepackage{textcomp} % provide euro and other symbols
\else % if luatex or xetex
  \usepackage{unicode-math} % this also loads fontspec
  \defaultfontfeatures{Scale=MatchLowercase}
  \defaultfontfeatures[\rmfamily]{Ligatures=TeX,Scale=1}
\fi
\usepackage{lmodern}
\ifPDFTeX\else
  % xetex/luatex font selection
\fi
% Use upquote if available, for straight quotes in verbatim environments
\IfFileExists{upquote.sty}{\usepackage{upquote}}{}
\IfFileExists{microtype.sty}{% use microtype if available
  \usepackage[]{microtype}
  \UseMicrotypeSet[protrusion]{basicmath} % disable protrusion for tt fonts
}{}
\makeatletter
\@ifundefined{KOMAClassName}{% if non-KOMA class
  \IfFileExists{parskip.sty}{%
    \usepackage{parskip}
  }{% else
    \setlength{\parindent}{0pt}
    \setlength{\parskip}{6pt plus 2pt minus 1pt}}
}{% if KOMA class
  \KOMAoptions{parskip=half}}
\makeatother
\usepackage{xcolor}
\usepackage[margin=1in]{geometry}
\setlength{\emergencystretch}{3em} % prevent overfull lines
\providecommand{\tightlist}{%
  \setlength{\itemsep}{0pt}\setlength{\parskip}{0pt}}
\setcounter{secnumdepth}{-\maxdimen} % remove section numbering
\ifLuaTeX
  \usepackage{selnolig}  % disable illegal ligatures
\fi
\usepackage{bookmark}
\IfFileExists{xurl.sty}{\usepackage{xurl}}{} % add URL line breaks if available
\urlstyle{same}
\hypersetup{
  hidelinks,
  pdfcreator={LaTeX via pandoc}}

\author{}
\date{}

\begin{document}

\section{Superlinearly many permanent values of prime-order binary
circulants}\label{superlinearly-many-permanent-values-of-prime-order-binary-circulants}

\textbf{Research manuscript, 20 November 2026. AI-assisted internal
mathematical checks have found no flaw in the argument below. External
review remains outstanding; this manuscript is not presented as
peer-reviewed.}

\subsection{Abstract}\label{abstract}

Let \(V(p)\) count the distinct permanents of all binary circulant
matrices of prime order \(p\). Nemeth, Seberry and Shu conjectured in
1979 that \(V(p)=O(p)\). We give an argument establishing \[
V(p)\gg \frac{p^{3/2}}{\sqrt{\log p}},
\] and hence contradicting that conjecture. The construction uses dense
binary circulants with prescribed Fourier profiles. A uniform permanent
estimate of McCullagh, in the form proved by Divoux and Wang, separates
permanent values within each row sum. The classical Brualdi--Newman
congruence separates values belonging to different row sums. The proof
does not establish the stronger conjectural fixed-row-sum asymptotic of
Resta and Sburlati. The probabilistic construction also admits a
deterministic, though exponential-time, specification.

\subsection{1. The question and the
result}\label{the-question-and-the-result}

For a positive integer \(n\) and a set
\(S\subseteq\mathbb Z/n\mathbb Z\), write \[
 C_n(S)_{ij}=\mathbf 1_{\{j-i\in S\}}.
\] Thus \(C_n(S)\) is a binary circulant, and every row and column has
sum \(|S|\). Define \[
 V(p)=\bigl|\{\operatorname{per}C_p(S):S\subseteq\mathbb Z/p\mathbb Z\}\bigr|.
\] All supports are allowed: zero may belong to \(S\), no symmetry
condition is imposed, and the row sum is not fixed.

Nemeth, Seberry and Shu {[}1, p.181{]} conjectured that \(V(p)=O(p)\)
for prime \(p\). This is Conjecture 27 in the Cheon--Wanless survey
{[}3{]}. Their discussion distinguishes it from the fixed-degree
conjecture of Resta and Sburlati {[}4{]}.

\textbf{Theorem 1.} There is an absolute constant \(c>0\) such that, for
every sufficiently large prime \(p\), \[
 V(p)\ge c\frac{p^{3/2}}{\sqrt{\log p}}.
\] In particular, \(V(p)/p\to\infty\) along the primes.

No explicit numerical starting prime is asserted. A less sharp, fully
specified amplitude grid also gives \[
 V(p)\ge \left(\frac13+o(1)\right)p^{4/3}.
\] These statements concern all row sums collectively, not a fixed
sparse row sum.

We use natural logarithms. For a real matrix,
\(\|\cdot\|_{\mathrm{op}}\) is the largest singular value and
\(\|\cdot\|_{\max}\) is the maximum absolute entry. The constants in all
asymptotic bounds below are independent of the support, row sum and
amplitude in their stated ranges.

\subsection{2. The external analytic
input}\label{the-external-analytic-input}

Put \(J_n=\mathbf1\mathbf1^T/n\).

\textbf{Analytic input.} Fix \(A<\infty\) and \(0\le r<1\). Uniformly
over doubly stochastic matrices \(D=J_n+B\) satisfying \[
 \|B\|_{\max}\le A/n,\qquad \|B\|_{\mathrm{op}}\le r,
\] one has \[
 \frac{n^n\operatorname{per}D}{n!}
 =\det(I-B^TB)^{-1/2}+O_{A,r}(n^{-1/2}).                 \tag{1}
\]

The qualitative version of (1) is McCullagh's permanent asymptotic
{[}5{]}, explicitly proved with uniformity in Divoux--Wang {[}6,
Proposition 3.1 and Remark 2.1{]}. The rate used here is extracted from
estimates in their proof; Appendix A gives that extraction. The
permanent asymptotic and its underlying analytic estimates are external
results, not claimed as contributions of this manuscript. In (1) the
determinant square root is the positive real one. Zero entries in \(D\)
are permitted.

For the qualitative conclusion \(V(p)/p\to\infty\), their stated uniform
\(o(1)\) estimate alone suffices. Section 6 separates that conclusion
from the quantitative extension, which uses the rate in (1).

\subsection{3. Binary supports with controlled Fourier
profiles}\label{binary-supports-with-controlled-fourier-profiles}

\textbf{Lemma 2.} Let \(p\ge5\) be prime, let \(k\) be an integer
satisfying \(p/3\le k\le2p/3\), and put \(\rho=k/p\). For each
\(a\in[0,1/12]\), there is a support \(S\subseteq\mathbb Z/p\mathbb Z\)
of size exactly \(k\) such that the eigenvalues of \(D=C_p(S)/k\)
satisfy \[
 \lambda_0=1,\qquad
 |\lambda_1-a/\rho|,|\lambda_{p-1}-a/\rho|\le\delta_p,
 \qquad
 \max_{\ell\notin\{0,1,p-1\}}|\lambda_\ell|\le\delta_p,       \tag{2}
\] where \(\delta_p=12\sqrt{\log p/p}\).

\textbf{Proof.} Write \(\zeta=e^{2\pi i/p}\). Independently select each
\(s\in\mathbb Z/p\mathbb Z\) with probability \[
 q_s=\rho+2a\cos(2\pi s/p).
\] These probabilities lie in \([1/6,5/6]\). If \(X_s\) is the selection
indicator, character orthogonality gives \[
 \mathbb E\sum_sX_s\zeta^{\ell s}
 =\begin{cases}
 k,&\ell=0,\\
 pa,&\ell=1,p-1,\\
 0,&\text{otherwise}.
 \end{cases}                                                   \tag{3}
\] For each frequency, Hoeffding's inequality applied to real and
imaginary parts bounds the probability that the complex error exceeds
\(t\) by \(4e^{-t^2/p}\). Taking \(t=2\sqrt{p\log p}\) and applying a
union bound over all \(p\) frequencies gives failure probability at most
\(4p^{-3}<1\). Thus a realization exists for which every Fourier error
is at most \(t\).

Its cardinality differs from \(k\) by at most \(t\), using frequency
zero. Add or delete that many elements to obtain a set \(S\) with
exactly \(k\) elements. Every edit changes every unnormalized Fourier
coefficient by at most one. All errors after repair are therefore at
most \(2t=4\sqrt{p\log p}\). Dividing by \(k\ge p/3\) proves (2). The
Fourier eigenvalue convention may reverse all frequencies; the estimates
are unchanged. \(\blacksquare\)

This is an existence argument for each profile. When many profiles are
used later, they may be chosen separately; no simultaneous random event
is required.

\subsection{4. The permanent associated with a
profile}\label{the-permanent-associated-with-a-profile}

Define \[
 h(x)=-\log(1-x)-x,
 \qquad
 F(\rho,a)=\frac12(\rho^{-1}-1)+h\bigl((a/\rho)^2\bigr).
\]

\textbf{Lemma 3.} There is an absolute constant \(C_0>0\) such that, for
all sufficiently large \(p\), every support supplied by Lemma 2
satisfies \[
 \left|\log\frac{\operatorname{per}C_p(S)}{k^p p!/p^p}
             -F(\rho,a)\right|
 \le C_0\sqrt{\frac{\log p}{p}}.                            \tag{4}
\]

\textbf{Proof.} A real circulant is normal and is diagonalized by the
Fourier matrix. For \(B=D-J_p\), the constant-direction eigenvalue is
zero and the other eigenvalues are \(\lambda_\ell\), \(\ell\ne0\).
Consequently the eigenvalues of \(B^TB\) are the squared magnitudes of
those eigenvalues.

Every entry of \(B\) is either \(-1/p\) or \(1/k-1/p\), so
\(\|B\|_{\max}\le2/p\). By (2), \[
 \|B\|_{\mathrm{op}}\le a/\rho+\delta_p\le1/4+\delta_p\le1/2
\] for all sufficiently large \(p\). We can apply (1) with the fixed
constants \(A=2\), \(r=1/2\).

Parseval's identity supplies the exact relation \[
 \sum_{\ell\ne0}|\lambda_\ell|^2=\frac pk-1=\rho^{-1}-1.       \tag{5}
\] Indeed, \(\sum_\ell|\sum_{s\in S}\zeta^{\ell s}|^2=pk\), and
\(\lambda_0=1\). The exceptional eigenvalues are complex conjugates. It
follows that \[
 \log\det(I-B^TB)^{-1/2}
 =\frac12(\rho^{-1}-1)+h(|\lambda_1|^2)
       +\frac12\sum_{\ell\notin\{0,1,p-1\}}h(|\lambda_\ell|^2). \tag{6}
\] For clarity, the last sum does not accumulate \(p\) copies of a
worst-case error. For \(0\le x<1\), \[
 0\le h(x)=\sum_{m\ge2}\frac{x^m}{m}
       \le\frac{x^2}{2(1-x)}.
\] Since the squared magnitudes in that sum are at most \(\delta_p^2\),
while their total is at most two by (5), their contribution is at most
\[
 \frac{\delta_p^2}{2(1-\delta_p^2)}=O(\log p/p).             \tag{7}
\] Take \(p\) large enough that \(\delta_p\le1/4\), and put
\(t=a/\rho\le1/4\). Then \[
 \bigl||\lambda_1|^2-t^2\bigr|
 \le\delta_p(2t+\delta_p)\le3\delta_p/4.
\] On \([0,1/4]\), \(h'(x)=x/(1-x)\le1/3\). The difference between the
exceptional term in (6) and \(h(t^2)\) is therefore at most
\(\delta_p/4\).

Finally, (1) changes the logarithm by \(O(p^{-1/2})\). This follows
directly because its determinant factor is at least one and its
normalized permanent is at least one by the van der Waerden inequality.
Combining these bounds with
\(\operatorname{per}C_p(S)=k^p\operatorname{per}D\) proves (4).
\(\blacksquare\)

\subsection{5. Separating different row
sums}\label{separating-different-row-sums}

The following congruence is classical: it is the binary case of
Brualdi--Newman {[}2, Theorem 3.2, p.161{]}. We include its proof.

\textbf{Lemma 4.} If \(p\) is prime, then \[
 \operatorname{per}C_p(S)\equiv |S|\pmod p.                 \tag{8}
\]

\textbf{Proof.} The permanent counts permutations \(\sigma\) of
\(\mathbb Z/p\mathbb Z\) for which \(\sigma(i)-i\in S\) for every \(i\).
Translation by \(c\) acts by conjugation, \[
 \sigma(i)\longmapsto\sigma(i-c)+c.
\] Every orbit has size one or \(p\). A fixed permutation commutes with
translation by one, so it is a translation \(\sigma(i)=i+s\). Such a
translation is admissible exactly when \(s\in S\). There are \(|S|\)
fixed permutations, proving (8). \(\blacksquare\)

Thus distinct degrees in the interval \([p/3,2p/3]\) produce disjoint
sets of permanent values.

\subsection{6. Counting distinct values}\label{counting-distinct-values}

\subsubsection{6.1. Qualitative disproof using only uniform
asymptotics}\label{qualitative-disproof-using-only-uniform-asymptotics}

Fix an arbitrary positive integer \(M\), and choose any \(M\) distinct
positive amplitudes in \([1/24,1/12]\). For each degree \(k\) in the
dense interval, construct a support for each amplitude.

For \(a>0\), \[
 \frac{\partial F}{\partial a}
 =\frac{2a^3}{\rho^4(1-a^2/\rho^2)}>0.                    \tag{9}
\] Continuity and compactness give a positive minimum separation of
adjacent limiting logarithmic values, uniformly over
\(\rho\in[1/3,2/3]\). Even with (4) replaced by a uniform \(o(1)\), all
\(M\) values within each degree are distinct for sufficiently large
\(p\). Lemma 4 separates the degrees. Therefore \[
 V(p)\ge M\bigl(\lfloor2p/3\rfloor-\lceil p/3\rceil+1\bigr).
\] Since \(M\) is arbitrary, \(V(p)/p\to\infty\). This already disproves
the original conjecture.

\subsubsection{6.2. Quantitative bound}\label{quantitative-bound}

On \(a\in[1/24,1/12]\), \(\rho\in[1/3,2/3]\), (9) gives \[
 \frac{\partial F}{\partial a}\ge\frac{3}{4096}.            \tag{10}
\] Set \(u_p=\sqrt{\log p/p}\), and choose \[
 M=\left\lfloor\frac{1}{131072 C_0u_p}\right\rfloor,
 \qquad a_j=\frac1{24}+\frac{j}{24M},\quad0\le j<M.
\] For sufficiently large \(p\), \(M\ge1\). Adjacent limiting
logarithmic values differ by at least \[
 \frac{1}{32768M}\ge4C_0u_p.
\] Each actual normalized logarithm differs from its target by at most
\(C_0u_p\), by Lemma 3. Hence all \(M\) permanents within each degree
are distinct. Lemma 4 again separates degrees, giving \[
 V(p)\ge M\bigl(\lfloor2p/3\rfloor-\lceil p/3\rceil+1\bigr)
       \gg p^{3/2}/\sqrt{\log p}.
\] This proves Theorem 1.

For a grid independent of the unspecified constant \(C_0\), take
\(M=\lfloor p^{1/3}\rfloor\) with the same amplitudes. The ratio of
adjacent target separation to the error grows like
\(p^{1/6}/\sqrt{\log p}\), and tends to infinity. This proves the stated
\((1/3+o(1))p^{4/3}\) corollary.

All constants in Lemmas 2 and 3 are independent of \(M\). The proof
therefore permits \(M\) to grow. Uniform qualitative convergence without
a rate would not, by itself, justify any prescribed growing choice of
\(M\).

\subsection{7. A deterministic
specification}\label{a-deterministic-specification}

This section gives a finite construction, not an efficient algorithm.
Let \(L=\lceil\log_2p\rceil\), and take a rational amplitude from the
grid with \(M=\lfloor p^{1/3}\rfloor\). Enumerate binary strings
lexicographically and select the first whose unnormalized Fourier
coefficients differ from the targets (3) by strictly less than
\(2\sqrt{pL}\), at every frequency. Such a string exists by the same
concentration argument, since \(L\ge\log p\) and the union failure
probability is less than one.

Repair its cardinality by adding or removing the smallest available
indices. The resulting error is less than
\(12\sqrt{L/p}=O(\sqrt{\log p/p})\), so the proof applies with enlarged
absolute constants. Squared-error comparisons involve real algebraic
numbers in a cyclotomic extension, and can be decided exactly. This
defines a deterministic exponential-time family; no closed-form support
or polynomial-time construction is claimed.

\subsection{Appendix A. Extracting a uniform
rate}\label{appendix-a.-extracting-a-uniform-rate}

We explain why the estimates in Divoux--Wang {[}6, proof of Proposition
3.8 and Section 3.3{]} imply the rate in (1). This appendix extracts a
rate from their proof; it does not reproduce their underlying Gaussian
and rook-polynomial analysis.

Fix \(A,r\), and choose a fixed compact interval \(I=[a,b]\) with
\(0<a<1<b<r^{-1}\). If \(r=0\), the estimate is immediate. For an
\(n\)-by-\(n\) matrix \(B\), define its rook polynomial by \[
 R_B(s)=\sum_{j=0}^n s^j
   \sum_{\substack{I,J\subseteq\{1,\ldots,n\}\\ |I|=|J|=j}}
     \operatorname{per}B[I,J],
\] where the permanent of the empty matrix is one. The normalized rook
expression in their proof is \[
 F_s(t)=R_B(st)\det(I-s^2t^2B^TB)^{1/2}
       =1+\sum_{k\ge1}f_k(s)t^k.
\] The determinant square root is the analytic branch equal to one at
\(t=0\); it is defined on the disk \(|st|r<1\). Their coefficient
estimates imply, uniformly for \(s\in I\), \[
 \sum_{1\le k\le K\sqrt n}|f_k(s)|
 \le\sum_{1\le E\le K\sqrt n}
       \frac{(CE/n)^{\lceil E/3\rceil}}{1-CE/n},           \tag{11}
\] with fixed constants \(C,K\). Their Cauchy estimate bounds the
remaining coefficient sum by \(O(e^{-c\sqrt n})\), after a fixed
sufficiently large choice of \(K\).

Here is a sharper evaluation of (11) than its qualitative use in that
paper. Put \(Q=\lceil K\sqrt n/3\rceil\) and \(A_1=3C\). For large
\(n\), the denominators are at least \(1/2\). Grouping by
\(q=\lceil E/3\rceil\), the sum is at most \[
 6\sum_{1\le q\le Q}(A_1q/n)^q.
\] Put \(v=A_1Q/n=O(n^{-1/2})\). For \(q\ge2\), \[
 (A_1q/n)^q\le(A_1/n)^2q^2v^{q-2}.
\] When \(v\le1/2\), the power sum on the right is at most 22. The
complete bound is thus \[
 6A_1/n+132A_1^2/n^2=O(n^{-1}).
\] Consequently their normalized rook approximation has uniform relative
error \(O(n^{-1})\) on \(I\).

Their Gamma representation is \[
 \frac{n^n\operatorname{per}D}{n!}=\mathbb E R_B(n/T),
 \qquad T\sim\operatorname{Gamma}(n+1,1).
\] Write \(Z=n/T\) and \(\Delta(s)=\det(I-s^2B^TB)^{-1/2}\). On \(I\),
\(\Delta\) and its derivative are uniformly bounded. Direct Gamma
moments give \[
 \mathbb EZ=1,\qquad\mathbb EZ^2=\frac n{n-1},\qquad
 \mathbb E|Z-1|\le(n-1)^{-1/2}.
\] Thus replacing \(\Delta(Z)\) by \(\Delta(1)\) on \(\{Z\in I\}\) costs
\(O(n^{-1/2})\). The probability of the complementary event is
exponentially small; the stronger weighted tail estimate in their
Section 3.3 bounds its contribution to \(\mathbb ER_B(Z)\) by
\(e^{-cn+O(\sqrt n)}\). Combining the central rook error, the Lipschitz
estimate and this tail bound proves (1), with constants depending only
on the fixed bounds.

\subsection{Scope, provenance and review
status}\label{scope-provenance-and-review-status}

The analytic permanent estimate and the modular congruence are prior
results. The proposed contribution is their application to a family of
binary Fourier profiles, followed by uniform amplitude separation and
counting across dense degrees. Literature searches checked the original
paper, subsequent surveys and specialist status listings; no prior full
resolution was located. Such searches cannot guarantee novelty or rule
out unindexed work.

Independent AI-assisted internal checks examined the argument, rate
extraction and conjecture scope. These checks are not independent human
peer review. Finite exact-permanent checks are supplied separately as
sanity checks, not as evidence sufficient for the asymptotic theorem. No
Lean verification or external peer review is claimed.

\subsection{References}\label{references}

\begin{enumerate}
\def\labelenumi{\arabic{enumi}.}
\tightlist
\item
  E. Nemeth, J. Seberry and M. Shu, \emph{On the distribution of the
  permanent of cyclic (0,1) matrices}, Utilitas Mathematica \textbf{16}
  (1979), 171--182; conjecture on printed p.181 (PDF p.11).
  \href{https://documents.uow.edu.au/~jennie/WEBPDF/071_1979.pdf}{Original
  paper}.
\item
  R. A. Brualdi and M. Newman, \emph{Some Theorems on the Permanent},
  Journal of Research of the National Bureau of Standards \textbf{69B}
  (1965), 159--163. Theorem 3.2.
  \href{https://nvlpubs.nist.gov/nistpubs/jres/69B/jresv69Bn3p159_A1b.pdf}{Original
  paper}.
\item
  G.-S. Cheon and I. M. Wanless, \emph{An update on Minc's survey of
  open problems involving permanents}, Linear Algebra and its
  Applications \textbf{403} (2005), 314--342. Conjecture 27,
  pp.329--330.
  \href{https://users.monash.edu.au/~iwanless/papers/permsurveyLAA.pdf}{Paper}.
\item
  G. Resta and G. Sburlati, \emph{On the number of different permanents
  of some sparse (0,1)-circulant matrices}, Linear Algebra and its
  Applications \textbf{375} (2003), 197--209.
  \href{https://doi.org/10.1016/S0024-3795(03)00649-9}{DOI}.
\item
  P. McCullagh, \emph{An asymptotic approximation for the permanent of a
  doubly stochastic matrix}, Journal of Statistical Computation and
  Simulation \textbf{84} (2014), 404--414.
  \href{https://arxiv.org/abs/1205.5723}{arXiv:1205.5723};
  \href{https://doi.org/10.1080/00949655.2012.712122}{DOI}.
\item
  A. Divoux and S. Wang, \emph{The Maximum Number of Hamiltonian Cycles
  in Tournaments}, arXiv:2610.23755v1 (2026). Proposition 3.1, Remark
  2.1, Proposition 3.8 and Section 3.3.
  \href{https://arxiv.org/html/2610.23755v1}{Paper}.
\end{enumerate}

The current specialist status page consulted during the novelty check
was
\href{https://yair-lavi.com/mincs-list/conjecture-27/index.html}{Yair
Lavi, Conjecture 27}. It is status evidence, not a source for the proof.

\end{document}
